\section*{第 \ref{ch:b2:amines} 章——胺与含氮化合物}

\begin{solution}{exo:b2:amines:1}
伯；仲；叔；季铵盐；伯（芳香）。
\end{solution}

\begin{solution}{exo:b2:amines:2}
乙酰胺 $<$ 苯胺 $<$ 氨 $<$ 甲胺（共轭酸的 $\mathrm pK_a$ 依次为
$-0.6$, 4.60, 9.26, 10.66).
\end{solution}

\begin{solution}{exo:b2:amines:3}
\textit{N}-甲基环己\omterm{def:b2:amines:imine}{亚胺}，\ce{C6H10=N-CH3}：胺与酮在弱酸性介质中
（pH 4--5），并移去生成的水（借助干燥剂或共沸蒸馏）。
\end{solution}

\begin{solution}{exo:b2:amines:4}
溴苯；苯甲腈；苯酚；4-羟基偶氮苯 \ce{C6H5-N=N-C6H4-OH}（偶联于苯氧负离子的
\ce{O-} 基团对位）。
\end{solution}

\begin{solution}{exo:b2:amines:5}
$\mathrm pK_a = 9.80$：pH 7 时，比值 $[\ce{(CH3)3NH+}]/[\ce{(CH3)3N}] = 10^{2.8} \approx 630$，
质子化形式占主导；柠檬汁（pH 接近 2）几乎将其全部质子化。振荡
胺的有机溶液与稀的酸水溶液，使胺以铵离子形式
转入水中；加碱则可使其重新游离出来。
\end{solution}

\begin{solution}{exo:b2:amines:6}
1-（环己-1-烯-1-基）吡咯烷。加成并脱水后，氮原子带
正电荷（成为亚铵离子），且无可失去的氢；转而从相邻碳原子上失去一个质子，
形成 \ce{C=C} 键。
\end{solution}

\begin{solution}{exo:b2:amines:7}
苯甲醛与丙-2-胺，或丙酮与苄胺，均在
pH 约为 6 时加入 \ce{NaBH3CN}。
\end{solution}

\begin{solution}{exo:b2:amines:8}
苯丙酮（1-苯基丙-1-酮），\ce{C6H5COCH2CH3}：一次加成生成\omterm{def:b2:amines:imine}{亚胺}阴离子，
水解得到酮。
\end{solution}

\begin{solution}{exo:b2:amines:9}
邻苯二甲酰\omterm{def:b2:amines:imine}{亚胺}钾 + 1-溴己烷：\textit{N}-己基邻苯二甲酰\omterm{def:b2:amines:imine}{亚胺}；再用肼（或水解）
使其释放出己-1-胺。被烷基化的邻苯二甲酰\omterm{def:b2:amines:imine}{亚胺}氮原子没有可用的孤电子对（它
离域于两个羰基上），且已无 \ce{N-H}，故不能再次烷基化。
\end{solution}

\begin{solution}{exo:b2:amines:10}
溴水使苯胺在氨基的邻位和对位发生溴代，
生成 2,4,6-三溴苯胺。\omterm{def:b2:amines:diazonium}{重氮化}后，再以 \ce{H3PO2} 将氨基换为氢：
此时三个溴原子彼此处于 1,3,5 位。
\end{solution}

\begin{solution}{exo:b2:amines:11}
4-硝基苯胺 + \ce{NaNO2} + \ce{HCl}，在 0--\qty{5}{\degreeCelsius} 下生成
4-硝基苯\omterm{def:b2:amines:diazonium}{重氮离子}氯化物。与苯酚在弱碱性溶液中偶联（苯氧负离子，活化作用强）
在 \ce{O-} 的对位发生，生成 4-[（4-硝基苯基）二氮烯基]苯酚，\ce{O2N-C6H4-N=N-C6H4-OH}，一种橙色
染料。
\end{solution}

\begin{solution}{exo:b2:amines:12}
三甲基铵基是较差而体积较大的离去基团，并使
\ce{CH3} 一侧的 $\beta$ 氢更具酸性、更易接近；碱夺取
取代程度较低的碳上的氢，生成 1-丁烯，即取代程度较低的烯烃（霍夫曼产物）。
\end{solution}

\begin{solution}{pb:b2:amines:1}
\textbf{1.} $\mathrm{{}^-O_3S{-}C_6H_4{-}NH_3^+}$：酸性的磺酸基已使胺质子化。
\textbf{2.} 两性离子在水中溶解度很小；碳酸根使铵基去质子化，
生成易溶的氨基苯磺酸钠。
\textbf{3.} \ce{NaNO2 + HCl -> HNO2 + NaCl}.
\textbf{4.}
\[
  \mathrm{{}^-O_3S{-}C_6H_4{-}NH_2 + HNO_2 + H^+ \to {}^-O_3S{-}C_6H_4{-}N_2^+ + 2\,H_2O} .
\]
\textbf{5.} 重氮盐受热时分解（生成苯酚与 \ce{N2}）；亚硝酸
同样会分解。
\textbf{6.} $5.00/173.19 = \qty{28.9}{mmol}$：需 \qty{28.9}{mmol} 的 \ce{NaNO2}，质量为 \qty{1.99}{g}。
\textbf{7.} 过量的亚硝酸会与偶联组分反应；可用淀粉–碘化钾
试纸检出，并以尿素或氨基磺酸将其破坏。
\textbf{8.} 其环受二甲氨基强烈活化，二甲氨基是强施主。
\textbf{9.} 位于 \ce{N(CH3)2} 的对位，属邻、对位定位，且邻位受
甲基位阻影响。
\textbf{10.} $\mathrm{{}^-O_3S{-}C_6H_4{-}N{=}N{-}C_6H_4{-}N(CH_3)_2}$.
\textbf{11.} 在强酸中二甲氨基被质子化：\ce{-NH(CH3)2+} 使环失活，
偶联随即停止。
\textbf{12.} 产物以质子化的、部分可溶的酸式存在；在碱性条件下，
磺酸钠在盐溶液中溶解度较小，沉淀析出。
\textbf{13.} 其正电荷离域到环上，末端氮是较差的
亲电中心；只有被 \ce{-O-}、\ce{-OH} 或 \ce{-NR2} 活化的环才能反应。
\textbf{14.} 碱性时呈黄橙色，酸性时呈红色。
\textbf{15.} 质子化发生在偶氮基的一个氮原子上（质子化形式通过与
二甲氨基的共轭而得到稳定）。
\textbf{16.} 约为 $\mathrm pK_a \pm 1$：大致在 3.1 至 4.4 之间，$\mathrm pK_a = 3.47$ 附近。
\textbf{17.} 两个苯环与偶氮基构成一个长共轭体系，最高
占据与最低空 $\pi$ 能级之间的能隙较小（见\cref{prop:b2:huckel:gap}），吸收
位于可见光区。
\textbf{18.} 质子化改变了 $\pi$ 体系中的电荷分布，并进一步缩小
能隙：吸收峰向长波方向移动，颜色由黄变红。
\textbf{19.} \qty{173.19}{g/mol}.
\textbf{20.} \qty{327.33}{g/mol}.
\textbf{21.} \qty{28.87}{mmol}.
\textbf{22.} $0.02887 \times 327.33 = \qty{9.45}{g}$.
\textbf{23.} 重氮盐分解、偶联不完全，以及产物溶于
母液和洗涤液中。
\textbf{24.} 甲基橙质量为 $0.80 \times 9.45 = \boldsymbol{\qty{7.56}{g}}$。
\end{solution}
